一块理想无吸收线性延迟片的两本征轴相对入射实验室 \(x\) 轴旋转角 \(\theta\)，本征轴之间相对相位延迟为 \(\delta\)。输入单色光为沿 \(x\) 方向的单位幅值线偏振 \(\mathbf E_{\rm in}=(1,0)^T\)，采用时间因子 \(e^{-i\omega t}\)，忽略共同相位。特别关注 \(\theta=\pi/4\) 和 \(\delta\) 接近 \(\pi/2\) 的情况。
(1) 给出实验室坐标 Jones 算符的旋转夹心表达式，验证无吸收条件。
(2) 求 \(\theta=\pi/4\) 的出射 Jones 矢量及 Stokes 分量，说明延迟符号与旋向的关联。
(3) 求产生圆偏振的条件以及半波片极限，并检查总强度不变。
(4) 若 \(\delta=\pi/2+\varepsilon\) 且 \(|\varepsilon|\ll1\)，定量求剩余线偏振和圆偏振度的首阶变化。
